\chapter{Development Workflow}
\label{chap:development}

This chapter gives the current build, boot, test, and code-reading workflow.

\section{Prerequisites}

\begin{itemize}
    \item \textbf{Rust nightly}, installed via \texttt{rustup}. The exact
        toolchain is pinned in the repository's
        \repofiletext|rust-toolchain.toml|{\texttt{rust-}\allowbreak\texttt{toolchain.toml}};
        do not copy a version
        from this manual.
    \item \textbf{QEMU} for AArch64 (\texttt{qemu-system-aarch64}) or x86-64
        (\texttt{qemu-system-x86\_64}). On macOS, TCG is the default and HVF
        is opt-in. Networking and network-backed services are enabled by
        default under TCG; the virtio-net smoke-test switch adds verification.
        HVF is unsuitable for the driver's direct EL0 device-MMIO path and
        therefore requires \texttt{--no-network}. KVM is an optional Linux
        accelerator, not a networking requirement.
        \textbf{HVF also does not honor hardware ASIDs} --- see the
        \repofiletext|docs/platforms/aarch64.md|{\emph{AArch64 Port Status}}
        write-up --- so HVF builds rely on the
        \texttt{hvf\_compat} whole-TLB-flush-on-switch fallback. Do not assume
        ASID-based TLB isolation works when debugging under \texttt{--hvf}.
        HVF also exposes no SMMU for protected NVMe DMA, so its compatibility
        boot runs the 15-test non-storage suite. Use TCG for the full 18-test
        suite, object-store tests, persistent Raft, and live upgrade.
    \item \textbf{Build dependencies:} \texttt{cargo}, \texttt{just} (optional,
        for convenience targets in the
        \repofiletext|Justfile|{\texttt{Justfile}}), and \texttt{mtools} for
        constructing x86-64 FAT boot images. VMware appliance packaging also
        requires \texttt{qemu-img}, normally installed with QEMU.
\end{itemize}

\section{Build targets}

CharlotteOS supports three architectures via custom target JSON specs in
\repodir|target_specs|:

\begin{longtable}{|p{5.8cm}|p{6.4cm}|}
\hline
\textbf{Target spec} & \textbf{Status} \\
\hline
AArch64 custom target & Original development target; all
    features and self-tests pass \\
x86-64 custom target & Four-LP ring-3 service, storage, IOMMU,
    discovery, distributed DNS, smoltcp TCP/IP, HTTP, deployment, migration,
    UTC time, and dynamic-membership paths pass under QEMU TCG; UEFI, VT-d,
    NVMe, and node-local services also pass under VMware Fusion \\
RISC-V custom target & Planned, not actively
    maintained \\
\hline
\end{longtable}

\subsection{Kernel build}

The commands below use
\repofiletext|target_specs/aarch64-unknown-none-catten.json|{the AArch64 target specification}
and
\repofiletext|target_specs/x86_64-unknown-none-catten.json|{the x86-64 target specification}.

\begin{verbatim}
# AArch64 (headless, serial console only):
cargo build --package catten \
  --target target_specs/aarch64-unknown-none-catten.json \
  --no-default-features --features acpi

# x86-64:
cargo build --package catten \
  --target target_specs/x86_64-unknown-none-catten.json \
  --no-default-features --features acpi
\end{verbatim}

The \texttt{display} feature requires a C library for the flanterm
framebuffer terminal. Headless boot uses serial output only (AArch64:
PL011 UART; x86-64: 16550 COM1).

\subsection{Userspace service programs}

Service binaries are built separately because they target \texttt{no\_std} /
\texttt{no\_main} with custom EL0 targets that depend on build-std:

\begin{verbatim}
# AArch64 EL0 service programs:
cargo build --release \
  --manifest-path crates/catten-services/Cargo.toml \
  --target crates/catten-services/aarch64-unknown-none.json \
  -Z build-std=core,alloc

# x86-64 ring-3 service programs:
cargo build --release \
  --manifest-path crates/catten-services/Cargo.toml \
  --target crates/catten-services/x86_64-unknown-none.json \
  -Z build-std=core,alloc

# Sitas user binary (requires external sitas crates):
cargo +nightly build --package catten-user \
  --target aarch64-unknown-none.json \
  -Zbuild-std=core,alloc
\end{verbatim}

\section{Booting in QEMU}

\subsection{AArch64}

The boot entry points are \repofile|scripts/boot-smp1.sh|,
\repofile|scripts/boot-smp2.sh|, and \repofile|scripts/run-aarch64.sh|.

Both architecture runners attach a NIC by default. Charlotte acquires a DHCP
lease and launches discovery, cluster, TCP/IP, HTTP, and time services without
a test option. The \texttt{*-test} switches only register additional verifiers;
\texttt{--no-network} is the explicit opt-out for an isolated boot.

\begin{verbatim}
# Single-LP self-test boot:
./scripts/boot-smp1.sh

# Two-LP self-test boot:
./scripts/boot-smp2.sh

# Or choose the profile, LP count, and timeout directly:
./scripts/run-aarch64.sh debug --smp 4 --timeout 180

# Apple Silicon: accelerated 15-test non-storage compatibility suite:
./scripts/run-aarch64.sh debug --hvf --no-network --smp 4 --timeout 180
\end{verbatim}

\subsection{x86-64}

The x86-64 runner is \repofile|scripts/run-x86_64.sh|;
\repofile|scripts/boot-x86.sh| is its single-instance convenience wrapper.
Both use mtools to construct the disposable FAT boot image, which avoids
platform-specific mounted-image discovery on macOS.

\begin{verbatim}
# Requires -cpu max for RDTSCP/FSGSBASE support:
./scripts/boot-x86.sh

# Select storage/IOMMU and run the four-LP bounded suite directly:
./scripts/run-x86_64.sh debug --smp 4 --block nvme --iommu intel \
  --fresh-storage --timeout 90

# Serve and host-probe the HTTP state report through smoltcp:
./scripts/run-x86_64.sh release --http-test --timeout 90

# In two terminals, run the complete four-LP cluster suite on both guests:
./scripts/run-x86_64.sh release --deploy-test --smp 4 \
  --instance node-a --mac 52:54:00:12:34:01 \
  --net-listen 12002 --timeout 360
./scripts/run-x86_64.sh release --deploy-test --smp 4 \
  --instance node-b --mac 52:54:00:12:34:02 \
  --net-connect 127.0.0.1:12002 --timeout 360
\end{verbatim}

Use \texttt{--tcpip-test} instead of \texttt{--deploy-test} for the smaller
two-guest smoltcp exchange. \texttt{--live-upgrade-test} runs the isolated
persistent service-manager handoff. The complete deployment suite is qualified
with four LPs; two LPs are not a supported test configuration for that load.

\section{Booting in VMware}

The appliance builder uses the same x86-64 build path without launching QEMU,
then converts the FAT boot image and a blank data image to VMware VMDKs:

\begin{verbatim}
./scripts/build-vmware-x86_64.sh release
\end{verbatim}

Open
\texttt{os-images/vmware/CharlotteOS.vmwarevm/CharlotteOS.vmx} in VMware
Fusion or Workstation. The SATA boot disk is nonpersistent; the NVMe data disk
is persistent. On first boot the guest formats that disk and installs the
missing signed x86-64 service artifacts. Later boots retain valid artifacts,
including deployed upgrades. The builder refuses to overwrite an existing
appliance unless \texttt{--replace} is supplied. The default 1~GiB data disk
requires the shared 4~MiB userspace heap so the object store can hold its
allocation bitmap, directory, and mirrored-directory mount buffer.

The supplied VM attaches VMware's emulated Intel 82574L adapter to NAT. The
userspace E1000E driver receives delegated MMIO, MSI-X, and protected-DMA
authority, waits for link without occupying an LP, and publishes the same
hardware-neutral \texttt{net0} service as virtio-net. Discovery, DNS, smoltcp,
HTTP, and cluster code therefore require no VMware-specific branch. The E1000E
path is qualified under Fusion, while paired VMware guests and VMXNET3 remain
unqualified. The complete setup and qualification record is in
\repofiletext|docs/platforms/vmware-x86_64.md|{the VMware x86-64 appliance guide}.

\section{Expected boot output}

\begin{verbatim}
CharlotteOS booting...
[init] BSP initialization complete.
[init] Secondary CPUs brought online.
[init] Starting scheduler...
[async] coordinator: submitted worker, awaiting completion...
[async] worker on LP1: sleeping 50ms
[async] worker on LP1: exiting
[async] COMPLETION RECEIVED, result Ok(42)
[async] SUCCESS: async syscall round-trip complete
Testing Complete. All Tests Passed!
\end{verbatim}

\section{Self-tests}

The repository combines host-side library tests with boot-time kernel and EL0
self-tests. Host suites cover components such as \texttt{catten-graft},
\texttt{charlotte-protocol-msg}, and \texttt{charlotte-smoltcp}. Boot tests run
from \texttt{self\_test::run\_self\_tests()} after initialization; assertion or
verifier failure fails the boot suite.

Deferred tests register an expected result bit and spawn their verifier through
\texttt{results::spawn\_verifier}. The result subsystem associates that
single verifier thread with the test in a fixed atomic table. If the verifier
panics, the panic handler uses a best-effort non-blocking scheduler lookup and
marks the associated test failed before emitting the panic. Consequently the
authoritative reporter names the failed test instead of leaving it pending
until the runner timeout. The panic path neither allocates nor waits for a
scheduler lock.

Representative boot suites include:

\begin{itemize}
    \item \repofiletext|crates/catten/src/self_test/completion.rs|{\texttt{completion.rs}} --- buffer ownership, cancel/deferred-reclaim,
        backpressure.
    \item \repofiletext|crates/catten/src/self_test/cq.rs|{\texttt{cq.rs}} --- CQ ring write/drain/overflow.
    \item \repofiletext|crates/catten/src/self_test/cq_completion.rs|{\texttt{cq\_completion.rs}} --- CQ ring integrated with completion
        subsystem.
    \item \repofiletext|crates/catten/src/self_test/syscall.rs|{\texttt{syscall.rs}} --- dispatch table routing via synthetic
        TrapFrame.
    \item \repofiletext|crates/catten/src/self_test/ipi.rs|{\texttt{ipi.rs}} --- bounded queue, backpressure, closure dispatch.
    \item \repofiletext|crates/catten/src/self_test/shard.rs|{\texttt{shard.rs}} --- \texttt{ShardLocal<T>} owner-check /
        re-entrancy guard, \texttt{ShardMailbox<M>} send/recv.
    \item \repofiletext|crates/catten/src/self_test/el0.rs|{\texttt{el0.rs}} --- real userspace thread (AArch64 EL0 and x86-64 ring 3).
    \item \repofiletext|crates/catten/src/self_test/el0_demo.rs|{\texttt{el0\_demo.rs}} --- EL0 userspace smoke test.
    \item \repofiletext|crates/catten/src/self_test/el0_ipc.rs|{\texttt{el0\_ipc.rs}} --- EL0 endpoint IPC between domains.
    \item \repofiletext|crates/catten/src/self_test/el0_net.rs|{\texttt{el0\_net.rs}} --- EL0 networking test.
    \item \repofiletext|crates/catten/src/self_test/el0_pingpong.rs|{\texttt{el0\_pingpong.rs}} --- bidirectional EL0 IPC.
    \item \repofiletext|crates/catten/src/self_test/el0_raft.rs|{\texttt{el0\_raft.rs}} --- two-node Raft election over local IPC.
    \item \repofiletext|crates/catten/src/self_test/el0_service.rs|{\texttt{el0\_service.rs}} --- EL0 service lifecycle (spawn, IPC,
        teardown).
    \item \repofiletext|crates/catten/src/self_test/el0_uart.rs|{\texttt{el0\_uart.rs}} --- userspace UART driver (delegated MMIO +
        IRQ, crash recovery).
    \item \repofiletext|crates/catten/src/self_test/el0_sitas.rs|{\texttt{el0\_sitas.rs}} --- sitas executor integration: sharded KV,
        the mailbox index demo, and the thread-join probe
        (\autoref{sec:sitas-mailbox-index}).
    \item \repofiletext|crates/catten/src/self_test/scheduler_lifecycle.rs|{\texttt{scheduler\_lifecycle.rs}} --- spawn, block, wake, abort,
        migration-safe migration, timer-wait affinity.
    \item \repofiletext|crates/catten/src/self_test/ipc.rs|{\texttt{ipc.rs}} --- IPC path tests (scalar, memory, delegation,
        reply tokens).
    \item \repofiletext|crates/catten/src/self_test/memory/allocator.rs|{\texttt{memory/allocator.rs}} --- heap allocator, stack allocator.
    \item \repofiletext|crates/catten/src/self_test/memory/object.rs|{\texttt{memory/object.rs}} --- memory object lifecycle tests.
    \item \repofiletext|crates/catten/src/self_test/adversarial.rs|{\texttt{adversarial.rs}} --- cancellation races, misuse paths,
        error-path testing.
\end{itemize}

\section{The async-syscall demo}

\repofile|crates/catten/src/demo.rs| is an end-to-end test of the
async model. It:

\begin{enumerate}
    \item Spawns a coordinator and worker from \texttt{bsp\_main}.
    \item The coordinator calls \texttt{submit\_worker}, which spawns a worker
        thread and returns a capability that fires when the worker exits.
    \item The worker sleeps 50ms (exercising the timer and block-yield-wake
        cycle).
    \item The worker exits; the exit-observer fires; the capability completes.
    \item The coordinator (blocked on \texttt{wait(cap)}) wakes and reads
        the result.
\end{enumerate}

This exercises thread creation, scheduling, timer events, the block--yield--wake
cycle, exit observation, and the wait/poll ABI. It is evidence for that path, not
a substitute for the wider self-test suite.

\section{CI pipeline}

GitHub Actions (\repofile|.github/workflows/ci.yml|):

\begin{enumerate}
    \item Checks formatting for the workspace and standalone EL0/tool crates.
    \item Runs warning-denied Clippy for the workspace, AArch64 services, and
        signing tool.
    \item Runs host-testable library suites and the signing self-test.
    \item Builds AArch64 and x86-64 kernels and the AArch64 service bundle.
    \item On pushes to \texttt{main}, boots the four-LP AArch64 QEMU suite and
        verifies its required success markers.
\end{enumerate}

The workflow runs on pushes and pull requests for its configured branches; the
QEMU boot job is restricted to pushes to \texttt{main}.

\section{Codebase reading guide}

For new contributors, the recommended reading order:

\begin{enumerate}
    \item \repofile|crates/catten/src/main.rs| --- kernel entry point, init
        sequence.
    \item \repofile|crates/catten/src/completion/mod.rs| --- the completion
        subsystem: capability table, submit/complete/poll/wait/cancel/close,
        \texttt{submit\_worker}, exit-observer.
    \item \repofile|crates/catten/src/completion/cq.rs| --- the CQ ring buffer.
    \item \repofile|crates/catten/src/syscall/mod.rs| --- syscall dispatch table.
    \item \repofile|crates/catten/src/demo.rs| --- the end-to-end async demo.
    \item \repodir|crates/catten/src/cpu/scheduler| --- threads, wakers,
        \texttt{RoundRobin}, \texttt{block\_thread}.
    \item \repodir|crates/catten/src/cpu/multiprocessor| --- SMP: IPI,
        \texttt{PerLp<T>}, \texttt{ShardLocal<T>}, shard mailboxes.
    \item \repodir|crates/catten/src/cpu/isa| --- ISA abstraction layer.
        Start with \repodirtext|crates/catten/src/cpu/isa/interface|{\texttt{interface/}}
        (traits), then
        \repodirtext|crates/catten/src/cpu/isa/aarch64|{\texttt{aarch64/}}
        (reference implementation).
    \item \repodir|crates/catten/src/service| --- supervisor, ELF loader,
        bootstrap.
    \item \repodir|crates/catten/src/device_management| --- device driver
        framework.
    \item \repodir|crates/catten-services/src/bin| --- reference EL0 services.
    \item \repofile|crates/catten-rt/src/lib.rs| --- userspace runtime.
\end{enumerate}

\section{Troubleshooting}

\begin{description}
    \item[Triple fault during boot] Check that the panic handler is working.
        Early faults before heap initialization rely on
        \texttt{early\_logln!} (direct serial writes). Ensure the serial port
        is configured correctly.
    \item[Blocked thread] The block--yield--wake cycle requires
        that \texttt{block\_thread} not clear \texttt{current\_handle} before
        \texttt{yield\_lp} saves the thread's context. Verify that the
        blocked thread is still the current thread at the point of the yield.
    \item[ASID-0 collision] The first user address space must not receive
        ASID 0 (which equals \texttt{KERNEL\_ASID}). The address-space table
        must be pre-populated with the kernel AS at index 0.
    \item[x86-64 LA57 mismatch] Derive the address width from
        \texttt{CR4.LA57} (the active paging mode), not from CPUID capability.
        Limine boots with 4-level paging even on CPUs that support 5-level.
\end{description}

\endinput
